\documentclass[runningheads]{llncs}

% ── Packages ──────────────────────────────────────────────────────────────────
\PassOptionsToPackage{table,dvipsnames}{xcolor}
\usepackage[review,year=2026,ID=230]{accv}
\usepackage{graphicx}
\usepackage{booktabs}
\usepackage[pagebackref,breaklinks,colorlinks,citecolor=blue]{hyperref}
\usepackage{amsmath,amssymb}
\usepackage{multirow}
\usepackage{array}
\usepackage{adjustbox}
\usepackage{enumitem}
\usepackage[section]{placeins}

\newcommand{\tabtight}{\setlength{\tabcolsep}{3.5pt}\renewcommand{\arraystretch}{1.12}}
\newcommand{\TDS}{\textsc{TDS}}

% ---- Revision marking (rebuttal round) ----------------------------------
% Mirrors main2.tex: added or rewritten text is blue; removed baseline text is
% suppressed because it remains available in supplementary.tex.
\usepackage{xcolor}
\usepackage[normalem]{ulem}  % normalem: do not hijack \emph
\definecolor{revblue}{RGB}{0,80,200}
\newcommand{\rev}[1]{\textcolor{revblue}{#1}}
\newcommand{\revdel}[1]{}
% Marks a section added during the rebuttal round. Beginning with S14, the
% document switches globally to revblue so body text and floats remain blue.
\newcommand{\revnote}{\noindent\textcolor{revblue}{\small\itshape
  [New in this revision, in response to the reviews.]}\par\vspace{2pt}}
\setlength{\emergencystretch}{4em}

\begin{document}
\raggedbottom

\title{\texorpdfstring{Supplementary Material:\\
  \rev{Measuring Spatiotemporal Sampling Sensitivity in Video Action Recognition:\\
  A Cross-Architecture Analysis at Scale}}{Supplementary Material: Measuring Spatiotemporal Sampling Sensitivity in Video Action Recognition: A Cross-Architecture Analysis at Scale}}

\author{Anonymous ACCV Submission}
\institute{}
\maketitle

\setcounter{section}{0}
\renewcommand\thesection{S\arabic{section}}
\renewcommand\thefigure{S\arabic{figure}}
\renewcommand\thetable{S\arabic{table}}
\setcounter{figure}{0} \setcounter{table}{0}

This document provides supplementary material for the main paper. We include: (S1) full coverage$\times$stride heatmaps for all 8 models and 8 datasets; (S2) Levene variance inflation analysis; (S3) full ANOVA effect size tables including FineGym; (S4) action sensitivity taxonomy detail; (S5) \rev{maximum-probability confidence-cascade curves}\revdel{entropy routing curves} for all models; (S6) clip duration analysis; (S7) spectral correlation detail; (S8) implementation details (splits, sampling, frame normalization, target-resolution fine-tuning, positional embedding interpolation); (S9) full multi-dataset spatial resolution results (all 8 datasets $\times$ 8 models); (S10) deterministic interactive dashboard and code availability; (S11) \TDS{} ranking robustness: leave-one-architecture-out, architecture-family ablation (CNN-only, Transformer-only, Transformer+SSM pools), and bootstrap confidence intervals; (S12) exploratory spectral and label-space evidence-localization checks showing why whole-frame flow is not a literal Nyquist proxy; and (S13) additional robustness checks on optimization headroom, coverage--stride interaction, and \rev{held-out cascade calibration}\revdel{entropy-threshold calibration}.

% ── S1: Full heatmaps ─────────────────────────────────────────────────────────
\section{Full Coverage$\times$Stride Heatmaps}
\label{sec:sup_heatmaps}

Figures~\ref{fig:heatmap_r3d}--\ref{fig:heatmap_videomamba} show the complete coverage$\times$stride accuracy heatmaps for all 8 architectures across all 8 datasets. Each cell represents Top-1 accuracy (\%) at a specific (coverage, stride) configuration. Color scale: green=high accuracy, red=low accuracy. FineGym heatmaps (all 8 models) are included; per-class optical-flow analysis for FineGym is reported in supplementary~S7.

Row order (coverage): 10\%, 25\%, 50\%, 75\%, 100\%. Column order (stride): 1, 2, 4, 8, 16.

\begin{figure}[p]
  \centering
  \includegraphics[width=\textwidth]{images/sup1_heatmap_r3d_18.pdf}
  \caption{R3D-18 (CNN-3D) coverage$\times$stride accuracy (\%) across 8 datasets. Sharp cliff at stride 4$\to$8 on AUTSL (75\%$\to$27\%) and similar cliff on FineGym are consistent with the high temporal demand of fine-grained actions. EPIC-Kitchens and UCF-101 retain high accuracy even at sparse sampling.}
  \label{fig:heatmap_r3d}
\end{figure}

\begin{figure}[p]
  \centering
  \includegraphics[width=\textwidth]{images/sup1_heatmap_mc3_18.pdf}
  \caption{MC3-18 (CNN-mix, channel-separated 3D) heatmaps. Similar pattern to R3D-18 but with slightly reduced \revdel{aliasing}\ \rev{evidence loss} on low-TDS datasets (UCF-101, DriveAct). FineGym shows the same severe cliff pattern as AUTSL.}
  \label{fig:heatmap_mc3}
\end{figure}

\begin{figure}[p]
  \centering
  \includegraphics[width=\textwidth]{images/sup1_heatmap_r2plus1d_18.pdf}
  \caption{R2+1D-18 (separable spatiotemporal convolution) heatmaps. Higher base accuracy than R3D-18 on several datasets (SSv2: 42.6\% vs 37.1\%), but a similar \revdel{aliasing-sensitivity}\ \rev{sampling-sensitivity} pattern.}
  \label{fig:heatmap_r2p1d}
\end{figure}

\begin{figure}[p]
  \centering
  \includegraphics[width=\textwidth]{images/sup1_heatmap_slowfast_r50.pdf}
  \caption{\rev{SlowFast-R50 (dual-pathway) heatmaps. The AUTSL cliff occurs at stride 2$\to$4 (77\%$\to$41\%), one step earlier than for the other CNNs under the original fixed-budget protocol. This difference is consistent with SlowFast's larger frame budget but does not isolate its causal role. On FineGym, SlowFast achieves the highest base accuracy (78.2\%) but suffers the largest evidence loss (71.5pp at stride=16) under this protocol.}}
  \label{fig:heatmap_slowfast}
\end{figure}

\begin{figure}[p]
  \centering
  \includegraphics[width=\textwidth]{images/sup1_heatmap_timesformer.pdf}
  \caption{\rev{TimeSformer (divided space--time attention) heatmaps. Accuracy is relatively flat across much of the original grid, with steeper degradation at the combined extreme of low coverage and high stride. FineGym drops 36.1pp at stride=16. These are descriptive fixed-budget curves; they do not identify the contribution of divided attention.}}
  \label{fig:heatmap_timesformer}
\end{figure}

\begin{figure}[p]
  \centering
  \includegraphics[width=\textwidth]{images/sup1_heatmap_vivit.pdf}
  \caption{\rev{ViViT (factorized attention) heatmaps. Despite being a Transformer, ViViT shows large fixed-budget drops (AUTSL: cliff at stride 2$\to$4, comparable to SlowFast); FineGym drops 55.2pp at stride=16. The comparison is observational and does not attribute this pattern to factorized temporal attention.}}
  \label{fig:heatmap_vivit}
\end{figure}

\begin{figure}[p]
  \centering
  \includegraphics[width=\textwidth]{images/sup1_heatmap_videomae.pdf}
  \caption{\rev{VideoMAE (masked autoencoder) heatmaps. Base accuracy is high on SSv2 (52.4\%), UCF-101 (95.4\%), and FineGym (78.0\%), while its original-protocol curve falls sharply on AUTSL (stride 4$\to$8) and FineGym (67.8pp at stride=16). This pattern cannot isolate the contribution of the MAE pretraining objective.}}
  \label{fig:heatmap_videomae}
\end{figure}

\begin{figure}[p]
  \centering
  \includegraphics[width=\textwidth]{images/sup1_heatmap_videomamba.pdf}
  \caption{VideoMamba (SSM selective scan) heatmaps (all 8 datasets). \revdel{On 5 of 8 datasets, VideoMamba shows a flat heatmap with minimal degradation across strides---the SSM recurrent state aggregates temporal context, providing natural robustness. FineGym (40.8pp), AUTSL (16.0pp), and SSv2 (14.1pp) are exceptions, revealing the limit of SSM temporal integration under high-frequency transitions.}\rev{Sensitivity is low on five datasets and higher on FineGym (40.8pp), AUTSL (16.0pp), and SSv2 (14.1pp); these observational differences do not isolate the effect of the temporal module.}}
  \label{fig:heatmap_videomamba}
\end{figure}

% ── S2: Levene variance ───────────────────────────────────────────────────────
\section{Levene Variance Inflation Analysis}
\label{sec:sup_levene}

Figure~\ref{fig:levene} plots inter-class standard deviation at stride=1 versus stride=16 for all model-dataset pairs. Points above the diagonal indicate that sparse sampling increases inter-class variability. Significant pairs ($p < 0.05$, Bonferroni-corrected) are highlighted in red.

\begin{figure}[!htbp]
  \centering
  \includegraphics[width=0.8\columnwidth]{images/sup2_levene_variance.pdf}
  \caption{Levene test: inter-class accuracy std at stride=1 vs.\ stride=16. Points above the dashed diagonal have higher variance under sparse sampling. Color encodes the variance ratio ($\sigma_{s16}/\sigma_{s1}$). \textbf{Key finding}: 67\% of model-dataset pairs show significant variance inflation ($p<0.05$), with the largest increases in VideoMAE/HMDB-51 (2.0$\times$) and R3D-18/HMDB-51 (1.85$\times$). TimeSformer and VideoMamba show stable variance ($\approx$1.1$\times$).}
  \label{fig:levene}
\end{figure}

Table~\ref{tab:levene_top} lists the 10 model-dataset pairs with highest variance inflation, quantifying the instability masked by mean accuracy.

\begin{table}[!htbp]
\centering
\caption{Top-10 variance inflation cases (stride=1$\to$16, coverage=100\%). All rows significant at $p{<}0.05$.}
\label{tab:levene_top}
\scriptsize\tabtight
\begin{tabular}{l l r r r r}
\toprule
Model & Dataset & $\sigma_{s=1}$ & $\sigma_{s=16}$ & Ratio & Levene $p$ \\
\midrule
VideoMAE    & HMDB-51      & 0.125 & 0.253 & 2.02$\times$ & $<$0.0001 \\
R3D-18      & HMDB-51      & 0.164 & 0.303 & 1.85$\times$ & 0.0001 \\
ViViT       & HMDB-51      & 0.158 & 0.277 & 1.75$\times$ & 0.0004 \\
SlowFast    & HMDB-51      & 0.187 & 0.326 & 1.74$\times$ & $<$0.0001 \\
SlowFast    & UCF-101      & 0.146 & 0.240 & 1.65$\times$ & $<$0.0001 \\
R2+1D       & HMDB-51      & 0.205 & 0.310 & 1.51$\times$ & 0.0021 \\
MC3-18      & HMDB-51      & 0.177 & 0.262 & 1.48$\times$ & 0.0187 \\
ViViT       & UCF-101      & 0.110 & 0.147 & 1.34$\times$ & 0.0093 \\
VideoMAE    & AUTSL        & 0.172 & 0.229 & 1.33$\times$ & 0.0324 \\
VideoMAE    & UCF-101      & 0.094 & 0.121 & 1.29$\times$ & 0.0239 \\
\bottomrule
\end{tabular}
\end{table}

% ── S3: Full ANOVA tables ─────────────────────────────────────────────────────
\section{\rev{Original Between-Cell ANOVA Effect Sizes}}
\label{sec:sup_anova}

\rev{This section preserves the originally submitted between-cell analysis for traceability. Because the same clips populate all 25 cells, its independence assumption is inappropriate; the corrected clip-level repeated-measures analysis in Sec.~\ref{sec:sup_rmanova} supersedes it for inference.}

Figure~\ref{fig:anova_bar} shows the mean $\eta^2$ per model (averaged over all valid datasets), with error bars indicating standard deviation.

\begin{figure}[!htbp]
  \centering
  \includegraphics[width=0.9\columnwidth]{images/sup3_anova_eta2.pdf}
  \caption{\rev{Originally submitted between-cell $\eta^2$ values. These descriptive bars are superseded by the repeated-measures partial $\eta^2$ estimates in Sec.~\ref{sec:sup_rmanova}.}\revdel{ANOVA effect sizes confirming architecture-dependent aliasing.}}
  \label{fig:anova_bar}
\end{figure}

Table~\ref{tab:anova_full} preserves the complete per-dataset values from the original between-cell model. \revdel{All stride effects are significant; coverage effects are significant in all cases.} \rev{We do not use those significance tests because they assume independent cells. The valid repeated-measures effect sizes and tests are reported in Sec.~\ref{sec:sup_rmanova}.}

\begin{table}[!htbp]
\centering
\caption{\rev{Original between-cell $\eta^2$ per model and dataset, retained for traceability only. Cov = $\eta^2(\text{coverage})$; Str = $\eta^2(\text{stride})$. See Sec.~\ref{sec:sup_rmanova} for corrected inference.}}
\label{tab:anova_full}
\scriptsize\tabtight
\begin{adjustbox}{max width=\textwidth}
\begin{tabular}{l *{8}{cc}}
\toprule
\multirow{2}{*}{Model} & \multicolumn{2}{c}{AUTSL} & \multicolumn{2}{c}{FineGym} & \multicolumn{2}{c}{SSv2} & \multicolumn{2}{c}{Diving-48} & \multicolumn{2}{c}{HMDB-51} & \multicolumn{2}{c}{DriveAct} & \multicolumn{2}{c}{EPIC}
 & \multicolumn{2}{c}{UCF-101} \\
\cmidrule(lr){2-3}\cmidrule(lr){4-5}\cmidrule(lr){6-7}\cmidrule(lr){8-9} \cmidrule(lr){10-11}\cmidrule(lr){12-13}\cmidrule(lr){14-15}\cmidrule(lr){16-17} & Cov & Str & Cov & Str & Cov & Str & Cov & Str & Cov & Str & Cov & Str & Cov & Str
 & Cov & Str \\
\midrule
VMamba & .88 & .09 & .77 & .16 & .88 & .09 & .94 & .04 & .90 & .09 & .83 & .15 & .95 & .05 & .88 & .07 \\
TSF    & .92 & .06 & .74 & .19 & .87 & .10 & .97 & .02 & .92 & .07 & .84 & .14 & .91 & .08 & .85 & .07 \\
MC3    & .66 & .18 & .50 & .33 & .54 & .28 & .81 & .13 & .81 & .18 & .75 & .22 & .89 & .11 & .79 & .18 \\
R3D    & .60 & .21 & .47 & .34 & .60 & .26 & .77 & .15 & .68 & .31 & .58 & .33 & .75 & .24 & .70 & .25 \\
R2+1D  & .63 & .20 & .50 & .32 & .62 & .26 & .76 & .17 & .72 & .27 & .49 & .38 & .80 & .18 & .76 & .20 \\
VMAE   & .69 & .18 & .53 & .28 & .69 & .22 & .82 & .12 & .84 & .16 & .43 & .43 & .87 & .11 & .81 & .14 \\
ViViT  & .61 & .24 & .53 & .32 & .51 & .33 & .69 & .20 & .71 & .27 & .60 & .31 & .80 & .19 & .68 & .26 \\
SF     & .51 & .29 & .26 & .50 & .44 & .36 & .51 & .33 & .55 & .39 & .39 & .46 & .67 & .30 & .61 & .36 \\
\bottomrule
\end{tabular}
\end{adjustbox}
\end{table}

% ── S4: Taxonomy ──────────────────────────────────────────────────────────────
\section{Action Sensitivity Taxonomy}
\label{sec:sup_taxonomy}

\begin{figure}[!htbp]
  \centering
  \includegraphics[width=\textwidth]{images/sup4_taxonomy.pdf}
  \caption{\revdel{Aliasing loss}\ \rev{Evidence loss} (pp) per sensitivity tier and dataset. Each bar shows mean accuracy drop (stride=1$\to$16) within the tier, annotated with the number of action classes. UCF-101 Low tier: $-0.3$pp; AUTSL remains high across tiers.}
  \label{fig:taxonomy_full}
\end{figure}

Table~\ref{tab:taxonomy_detail} lists the threshold values that define tier boundaries per dataset, derived from the 33rd and 67th percentiles of the per-class accuracy drop.

\begin{table}[!htbp]
\centering
\caption{Taxonomy tier thresholds (accuracy drop, pp): classes below the Low threshold alias negligibly; classes above the High threshold are severely affected. Thresholds are dataset-specific and computed across all valid architectures.}
\label{tab:taxonomy_detail}
\scriptsize\tabtight
\begin{tabular}{l r r r r r}
\toprule
Dataset & Classes & Low threshold & High threshold & High count & Low count \\
\midrule
AUTSL         & 226 & 66.9pp & 73.9pp & 75 & 74 \\
Diving-48     &  47 & 25.0pp & 47.7pp & 16 & 16 \\
SSv2          & 174 & 59.7pp & 76.9pp & 58 & 57 \\
HMDB-51       &  51 &  8.4pp & 28.5pp & 17 & 17 \\
EPIC-Kitchens &  89 &  0.0pp & 19.4pp & 30 & 11 \\
DriveAct      &  33 & 14.8pp & 51.0pp & 11 & 11 \\
UCF-101       & 101 &  1.4pp &  7.5pp & 35 & 31 \\
\bottomrule
\end{tabular}
\end{table}

% ── S5: Routing curves ────────────────────────────────────────────────────────
\section{\rev{Maximum-Probability Confidence-Cascade Curves --- All Models}}
\label{sec:sup_routing}

\revdel{Figure~\ref{fig:routing_all} shows accuracy vs. average frames for all 8 architectures on SSv2 under the entropy routing protocol.} \rev{Figure~\ref{fig:routing_all} shows the descriptive in-sample accuracy--frame curves obtained by thresholding maximum class probability. These curves are not held-out estimates; held-out calibration results are reported in Sec.~\ref{sec:sup_routing_heldout}.} Each curve sweeps $\tau \in [0.10, 0.95]$. Dashed horizontal lines indicate fixed-budget baselines (4f, 16f) and the oracle.

\begin{figure}[!htbp]
  \centering
  \includegraphics[width=\textwidth]{images/sup5_routing_all_models.pdf}
  \caption{\rev{Maximum-probability cascade: descriptive in-sample accuracy vs. average frames on SSv2 (4f cheap $\to$ 16f dense). Held-out threshold selection is analysed separately in Sec.~\ref{sec:sup_routing_heldout}.}\revdel{Entropy routing: accuracy vs. average frames on SSv2.}}
  \label{fig:routing_all}
\end{figure}

Table~\ref{tab:routing_full} provides the full \rev{in-sample cascade frontier}\revdel{entropy routing results} across all models and datasets at the optimal threshold (avg frames $\leq 8$); it is retained as a descriptive diagnostic, not an unbiased performance estimate.

\begin{table}[!htbp]
\centering
\caption{\rev{In-sample maximum-probability cascade frontier: best Top-1 accuracy at avg $\leq8$ frames per model/dataset. Held-out results and uncertainty are in Sec.~\ref{sec:sup_routing_heldout}.}\revdel{Entropy routing: best accuracy at avg $\leq8$ frames per model/dataset.}}
\label{tab:routing_full}
\scriptsize\tabtight
\begin{adjustbox}{max width=\textwidth}
\begin{tabular}{l cccccccc r}
\toprule
Model & AUTSL & FineGym & Diving & SSv2 & HMDB & DriveAct & EPIC & UCF & Avg \\
\midrule
R3D-18       & 73.3 & 68.5 & 27.6 & 34.6 & 80.2 & 66.1 & 36.7 & 81.6 & 58.6 \\
MC3-18       & 63.4 & 62.9 & 31.8 & 32.4 & 78.5 & 69.4 & 36.5 & 85.6 & 57.6 \\
R2+1D        & 74.2 & 72.9 & 36.3 & 40.8 & 73.3 & 59.6 & 34.8 & 89.0 & 60.1 \\
SlowFast     & 59.5 & 57.8 & 46.6 & 35.0 & 76.3 & 67.4 & 37.4 & 87.9 & 58.5 \\
TimeSformer  & 67.1 & 66.8 & 38.9 & 42.5 & 80.6 & 68.5 & 32.4 & 91.3 & \textbf{61.0} \\
ViViT        & 67.2 & 66.1 & 48.9 & 32.2 & 80.1 & 64.3 & 31.3 & 94.3 & 60.6 \\
VideoMAE     & \textbf{79.5} & \textbf{76.4} & \textbf{52.6} & \textbf{51.9} & \textbf{84.7} & \textbf{69.0} & \textbf{37.8} & \textbf{95.9} & \textbf{68.5} \\
VideoMamba   & 65.0 & 73.2 & 37.4 & 44.4 & 69.8 & 57.6 & 28.4 & 88.5 & 58.0 \\
\bottomrule
\end{tabular}
\end{adjustbox}
\end{table}

% ── S6: Clip duration ─────────────────────────────────────────────────────────
\section{Clip Duration Analysis}
\label{sec:sup_duration}

\begin{figure}[!htbp]
  \centering
  \includegraphics[width=\textwidth]{images/sup6_clip_duration.pdf}
  \caption{\revdel{Aliasing loss}\ \rev{Evidence loss} (pp) vs.\ clip duration. \rev{Shorter clips lose more accuracy} (Pearson $r=-0.3$ to $-0.8$ across models), consistent with lower temporal redundancy. Diving-48 is the exception.}
  \label{fig:duration}
\end{figure}

% ── S7: Spectral correlation ──────────────────────────────────────────────────
\section{\rev{Exploratory Motion-Proxy Analysis}}
\label{sec:sup_spectral}

We estimated per-class temporal activity using optical-flow magnitude (Farneback algorithm~\cite{farneback2003two}) computed on 5 representative videos per class, sampling 16 frames per video at 112$\times$112px. \rev{This is an exploratory motion proxy, not a direct measurement of action frequency or a validation of classical Nyquist aliasing.}

Table~\ref{tab:spectral} reports Pearson correlation between per-class mean flow magnitude and \rev{evidence loss}\revdel{aliasing loss} (accuracy drop stride=1$\to$16). The range $r=-0.475$ to $0.327$ shows that whole-frame motion is an inconsistent proxy: camera motion can dominate flow, localized hand motion can be diluted, and a brief discriminative event need not have high mean motion.

\begin{table}[!htbp]
\centering
\caption{\rev{Pearson $r$ between whole-frame optical-flow magnitude and per-class evidence loss. The mixed signs show that flow magnitude is not a literal Nyquist proxy.}\revdel{Positive $r$ supports the Nyquist hypothesis.}}
\label{tab:spectral}
\scriptsize\tabtight
\begin{tabular}{l r r r}
\toprule
Dataset & Pearson $r$ & $p$-value & Significant ($p{<}0.05$) \\
\midrule
DriveAct      & 0.327 & 0.065 & Marginal \\
AUTSL         & 0.285 & $<$0.001 & Yes \\
SSv2          & 0.184 & $<$0.001 & Yes \\
HMDB-51       & 0.180 & 0.206 & No \\
Diving-48     & 0.140 & 0.349 & No \\
UCF-101       & 0.031 & 0.762 & No \\
EPIC-Kitchens & \rev{$-0.002$}\revdel{--0.002} & 0.985 & No \\
FineGym       & \rev{$-0.475$}\revdel{--0.475} & $<$0.0001 & Yes (inverted) \\
\bottomrule
\end{tabular}
\end{table}

% ── S8: Implementation details ────────────────────────────────────────────────
\section{Implementation Details}
\label{sec:sup_impl}

\paragraph{Hardware.}
Training experiments used NVIDIA A100 (40GB) and H200 NVL (141GB) GPUs (CUDA 13.2). Temporal-sweep evaluation (1,600 configurations) was parallelized across 8 nodes using a custom Slurm auto-submitter. Inference latency measurements used an NVIDIA RTX PRO 6000 Blackwell GPU (96GB) with bfloat16 autocast, batch size=1, 20 warmup iterations, 100 benchmark iterations, and CUDA event timing with explicit synchronization.

\paragraph{Splits, clip sampling, and evaluation determinism.}
We use the official train/validation/test split when a dataset provides one. When the public split is larger than needed for the controlled sweep, we form a deterministic balanced evaluation subset with approximately 20 validation clips per class; if a class has fewer available clips, all available validation clips are used. The same clip list is reused across all architectures, resolutions, coverages, and strides for a dataset. This makes pairwise stride/resolution differences comparable and avoids model-specific sampling noise. Reported main-paper numbers are validation-set measurements; no test-set tuning is used.

\paragraph{Seeds and uncertainty.}
Because the full grid requires thousands of model--dataset--resolution evaluations, each model--dataset fine-tuning run uses one fixed training seed and deterministic evaluation. Confidence intervals in S11 are therefore not random-seed intervals: they bootstrap over the 8-architecture pool. Per-class taxonomy and Levene analyses use class-level variation, and routing curves use validation-set threshold sweeps rather than an independently calibrated test threshold.

\paragraph{Frame and crop normalization.}
Training uses random resized crop, horizontal flip, and ImageNet normalization. Evaluation uses deterministic resize/center-crop at the requested spatial resolution. For every coverage/stride setting, frames are selected from the same leading candidate window. Each architecture receives its expected input length; if a candidate window is too short, we pad by repeating the last available frame, so stride changes alter sampling density rather than tensor shape.

\paragraph{Training hyperparameters.}
All models fine-tuned for 10 epochs from Kinetics-400 checkpoints. Transformers and VideoMamba: AdamW, lr=$2\times10^{-5}$, weight decay=0.05, batch size=32 per GPU. CNNs: AdamW, lr=$10^{-4}$, batch size=128 per GPU. SlowFast: AdamW, lr=$10^{-4}$, batch size=32 per GPU. All models use cosine learning rate decay with linear warmup (1 epoch). Spatial augmentation: random resized crop, horizontal flip, normalization (ImageNet statistics).

\paragraph{Positional embedding interpolation for Transformers and SSMs.}
For target-resolution fine-tuning at non-native resolutions, Transformer and SSM positional embeddings must be adapted. \textit{The wrong approach}: using the HuggingFace mismatched-size flag discards positional embeddings entirely, causing a large accuracy cliff (typically 10--20pp worse than our approach). \textit{Our approach}: load the native-resolution checkpoint, then bicubically interpolate spatial positional embeddings from the native grid ($14{\times}14{=}196$ tokens at 224px) to the target grid ($\lfloor r/16\rfloor^2$ tokens at resolution $r$) before copying weights to a model instantiated at the target resolution. This preserves the spatial topology learned during pretraining and ensures fine-tuning starts from a meaningful positional encoding rather than a random one. The same bicubic interpolation applies to TimeSformer, VideoMAE, ViViT, and VideoMamba.

\paragraph{SlowFast spatial pooling fix.}
The original SlowFast-R50 head uses a fixed $7\times7$ average pooling kernel, assuming 224px input. For non-native resolutions, we replace both head pooling layers with \texttt{AdaptiveAvgPool3d(1)}, enabling forward passes at arbitrary input resolutions without parameter changes.

\paragraph{FineGym dataset details.}
FineGym~\cite{shao2020finegym} contains 97 fine-grained gymnastics action classes (floor exercise, uneven bars, balance beam, vault; parallel bars, horizontal bar, pommel horse, rings for men). We use the FineGym288 subset (1,960 clips), with deterministic balanced clip lists for the controlled sweeps. Preprocessing: clips are trimmed to the annotated action boundary (mean duration 3.2s at 25fps) and padded with the last frame if shorter than 16 frames.

% ── S9: Multi-dataset spatial resolution ──────────────────────────────────────
\section{Multi-Dataset Spatial Resolution Results}
\label{sec:sup_spatial}

Tables~\ref{tab:sup_sp_ssv2}--\ref{tab:sup_sp_finegym} extend the spatial robustness analysis from the main paper (Section~4.3) to all 8 datasets. Tables \ref{tab:sup_sp_ssv2}--\ref{tab:sup_sp_epic} show spatial-resolution results at 96--224px. Table~\ref{tab:sup_sp_finegym} shows FineGym target-resolution fine-tuning results at 48--224px (one fine-tuned checkpoint per resolution). \textbf{Bold} = native training resolution. $\Delta_\text{max}$ = largest accuracy drop from native across all tested resolutions.

\rev{For the no-fine-tuning tables, ``resolution'' denotes the square pre-resize applied to decoded source frames before they are passed through the unchanged native checkpoint and its stored preprocessing. These are therefore input-scale/resampling stress tests, not evaluations of a Transformer instantiated on a new positional-token grid. Positional embedding interpolation is used only in the target-resolution fine-tuning branch.}

\begin{table}[!htbp]
\centering
\caption{\rev{Input-scale sensitivity on \textbf{SSv2} under source-frame pre-resizing and unchanged native checkpoint preprocessing (repeated from main paper; no fine-tuning).}}
\label{tab:sup_sp_ssv2}
\scriptsize\tabtight
\begin{adjustbox}{max width=\textwidth}
\begin{tabular}{l r rrrr r}
\toprule
Model & Native & 96px & 112px & 160px & 224px & $\Delta_\text{max}$ \\
\midrule
R3D-18      & 112px & 35.9 & \textbf{37.1} & 30.3 & 17.2 & $-19.9$pp \\
MC3-18      & 112px & 32.4 & \textbf{33.6} & 27.8 & 14.8 & $-18.8$pp \\
R2+1D       & 112px & 40.1 & \textbf{42.6} & 36.2 & 20.8 & $-21.8$pp \\
SlowFast    & 224px & 27.7 & 32.4 & 45.7 & \textbf{49.5} & $-21.8$pp \\
\midrule
TimeSformer & 224px & 39.3 & 39.4 & 41.4 & \textbf{42.3} & $-3.0$pp \\
ViViT       & 224px & 36.4 & 37.1 & 38.1 & \textbf{38.3} & $-1.9$pp \\
VideoMAE    & 224px & 48.9 & 49.2 & 51.5 & \textbf{52.3} & $-3.4$pp \\
VideoMamba  & 224px & 37.5 & 39.3 & 42.5 & \textbf{43.9} & $-6.4$pp \\
\bottomrule
\end{tabular}
\end{adjustbox}
\end{table}

\begin{table}[!htbp]
\centering
\caption{Spatial robustness on \textbf{UCF-101} (no retraining).}
\label{tab:sup_sp_ucf}
\scriptsize\tabtight
\begin{adjustbox}{max width=\textwidth}
\begin{tabular}{l r rrrr r}
\toprule
Model & Native & 96px & 112px & 160px & 224px & $\Delta_\text{max}$ \\
\midrule
R3D-18      & 112px & 79.9 & \textbf{81.2} & 75.3 & 55.1 & $-26.1$pp \\
MC3-18      & 112px & 82.9 & \textbf{85.4} & 81.0 & 63.9 & $-21.5$pp \\
R2+1D       & 112px & 88.0 & \textbf{88.6} & 85.1 & 64.9 & $-23.7$pp \\
SlowFast    & 224px & 69.7 & 75.8 & 87.5 & \textbf{91.4} & $-21.7$pp \\
\midrule
TimeSformer & 224px & 89.1 & 89.7 & 90.7 & \textbf{91.5} & $-2.4$pp \\
ViViT       & 224px & 92.0 & 93.5 & 94.7 & \textbf{95.1} & $-3.1$pp \\
VideoMAE    & 224px & 95.3 & 95.7 & 96.1 & \textbf{96.4} & $-1.1$pp \\
VideoMamba  & 224px & 76.7 & 78.7 & 82.6 & \textbf{83.5} & $-6.8$pp \\
\bottomrule
\end{tabular}
\end{adjustbox}
\end{table}

\begin{table}[!htbp]
\centering
\caption{Spatial robustness on \textbf{HMDB-51} (no retraining).}
\label{tab:sup_sp_hmdb}
\scriptsize\tabtight
\begin{adjustbox}{max width=\textwidth}
\begin{tabular}{l r rrrr r}
\toprule
Model & Native & 96px & 112px & 160px & 224px & $\Delta_\text{max}$ \\
\midrule
R3D-18      & 112px & 78.8 & \textbf{80.3} & 75.9 & 59.8 & $-20.5$pp \\
MC3-18      & 112px & 76.7 & \textbf{78.6} & 74.2 & 51.0 & $-27.6$pp \\
R2+1D       & 112px & 74.4 & \textbf{73.1} & 68.4 & 50.5 & $-22.6$pp \\
SlowFast    & 224px & 59.5 & 67.2 & 80.5 & \textbf{81.6} & $-22.1$pp \\
\midrule
TimeSformer & 224px & 75.9 & 77.2 & 77.5 & \textbf{78.9} & $-3.0$pp \\
ViViT       & 224px & 78.0 & 79.0 & 79.9 & \textbf{79.5} & $-1.5$pp \\
VideoMAE    & 224px & 83.7 & 83.1 & 84.0 & \textbf{83.8} & $-0.7$pp \\
VideoMamba  & 224px & 59.9 & 63.8 & 67.2 & \textbf{69.8} & $-9.9$pp \\
\bottomrule
\end{tabular}
\end{adjustbox}
\end{table}

\begin{table}[!htbp]
\centering
\caption{Spatial robustness on \textbf{Diving-48} (no retraining).}
\label{tab:sup_sp_div}
\scriptsize\tabtight
\begin{adjustbox}{max width=\textwidth}
\begin{tabular}{l r rrrr r}
\toprule
Model & Native & 96px & 112px & 160px & 224px & $\Delta_\text{max}$ \\
\midrule
R3D-18      & 112px & 26.5 & \textbf{28.8} & 22.1 & 12.8 & $-16.0$pp \\
MC3-18      & 112px & 30.8 & \textbf{31.6} & 27.3 & 13.7 & $-17.9$pp \\
R2+1D       & 112px & 33.8 & \textbf{35.3} & 30.1 & 18.7 & $-16.6$pp \\
SlowFast    & 224px & 18.4 & 25.0 & 43.3 & \textbf{50.5} & $-32.1$pp \\
\midrule
TimeSformer & 224px & 37.4 & 39.3 & 38.6 & \textbf{38.1} & $-0.7$pp \\
ViViT       & 224px & 48.0 & 49.9 & 52.8 & \textbf{53.0} & $-5.0$pp \\
VideoMAE    & 224px & 48.7 & 48.3 & 50.9 & \textbf{48.6} & $-0.3$pp \\
VideoMamba  & 224px & 29.2 & 29.1 & 36.1 & \textbf{36.6} & $-7.5$pp \\
\bottomrule
\end{tabular}
\end{adjustbox}
\end{table}

\begin{table}[!htbp]
\centering
\caption{Spatial resolution results on \textbf{AUTSL} (sign language). CNN degradation is severe under off-native evaluation; SlowFast drops from 79.2\% at 224px to 22.5\% at 96px ($-56.7$pp).}
\label{tab:sup_sp_autsl}
\scriptsize\tabtight
\begin{adjustbox}{max width=\textwidth}
\begin{tabular}{l r rrrr r}
\toprule
Model & Native & 96px & 112px & 160px & 224px & $\Delta_\text{max}$ \\
\midrule
R3D-18      & 112px & 72.2 & \textbf{75.0} & 63.0 & 27.9 & $-47.1$pp \\
MC3-18      & 112px & 59.1 & \textbf{63.7} & 52.7 & 12.2 & $-51.5$pp \\
R2+1D       & 112px & 73.7 & \textbf{75.9} & 68.9 & 38.4 & $-37.5$pp \\
SlowFast    & 224px & 22.5 & 32.6 & 70.0 & \textbf{79.2} & $-56.7$pp \\
\midrule
TimeSformer & 224px & 65.1 & 66.0 & 68.3 & \textbf{68.5} & $-3.4$pp \\
ViViT       & 224px & 49.2 & 51.8 & 53.3 & \textbf{53.6} & $-4.4$pp \\
VideoMAE    & 224px & 77.3 & 78.2 & 79.5 & \textbf{79.6} & $-2.3$pp \\
VideoMamba  & 224px & 30.4 & 34.2 & 53.3 & \textbf{65.8} & $-35.4$pp \\
\bottomrule
\end{tabular}
\end{adjustbox}
\end{table}

\begin{table}[!htbp]
\centering
\caption{Spatial robustness on \textbf{DriveAct} (in-cabin activities; no retraining).}
\label{tab:sup_sp_drive}
\scriptsize\tabtight
\begin{adjustbox}{max width=\textwidth}
\begin{tabular}{l r rrrr r}
\toprule
Model & Native & 96px & 112px & 160px & 224px & $\Delta_\text{max}$ \\
\midrule
R3D-18      & 112px & 65.0 & \textbf{68.3} & 52.7 & 19.0 & $-49.3$pp \\
MC3-18      & 112px & 65.2 & \textbf{69.0} & 38.4 &  6.2 & $-62.8$pp \\
R2+1D       & 112px & 67.4 & \textbf{62.5} & 40.0 & 12.9 & $-49.6$pp \\
SlowFast    & 224px & 31.7 & 40.2 & 65.8 & \textbf{70.1} & $-38.4$pp \\
\midrule
TimeSformer & 224px & 68.3 & 69.2 & 69.6 & \textbf{70.3} & $-2.0$pp \\
ViViT       & 224px & 59.4 & 58.7 & 62.1 & \textbf{63.4} & $-4.7$pp \\
VideoMAE    & 224px & 71.4 & 69.9 & 71.7 & \textbf{73.0} & $-3.1$pp \\
VideoMamba  & 224px & 35.0 & 38.2 & 46.9 & \textbf{57.8} & $-22.8$pp \\
\bottomrule
\end{tabular}
\end{adjustbox}
\end{table}

\begin{table}[!htbp]
\centering
\caption{Spatial robustness on \textbf{EPIC-Kitchens} (egocentric cooking; no retraining).}
\label{tab:sup_sp_epic}
\scriptsize\tabtight
\begin{adjustbox}{max width=\textwidth}
\begin{tabular}{l r rrrr r}
\toprule
Model & Native & 96px & 112px & 160px & 224px & $\Delta_\text{max}$ \\
\midrule
R3D-18      & 112px & 33.3 & \textbf{37.2} & 26.7 &  8.6 & $-28.6$pp \\
MC3-18      & 112px & 36.1 & \textbf{36.2} & 27.2 &  9.4 & $-26.8$pp \\
R2+1D       & 112px & 32.8 & \textbf{35.5} & 27.6 & 11.8 & $-23.7$pp \\
SlowFast    & 224px & 11.3 & 17.7 & 31.3 & \textbf{37.5} & $-26.2$pp \\
\midrule
TimeSformer & 224px & 20.5 & 19.2 & 25.6 & \textbf{31.7} & $-12.5$pp \\
ViViT       & 224px & 30.2 & 34.1 & 36.1 & \textbf{37.9} & $-7.7$pp \\
VideoMAE    & 224px & 30.9 & 32.6 & 36.5 & \textbf{38.6} & $-7.7$pp \\
VideoMamba  & 224px & 11.1 & 12.6 & 20.8 & \textbf{28.3} & $-17.2$pp \\
\bottomrule
\end{tabular}
\end{adjustbox}
\end{table}

\begin{table}[!htbp]
\centering
\caption{\textbf{FineGym} spatial resolution results after target-resolution fine-tuning (one checkpoint per resolution, 10 fine-tuning epochs from native checkpoint). Bold = native training resolution (112px for CNNs, 224px for Transformers/SSM). Unlike other datasets in S9 (cross-resolution without retraining), FineGym shows CNN performance that \emph{improves beyond native} after fine-tuning at higher resolutions (R3D-18: 71.5\%@112px$\to$76.1\%@224px), consistent with higher spatial resolution aiding gymnastics pose discrimination.}
\label{tab:sup_sp_finegym}
\scriptsize\tabtight
\begin{adjustbox}{max width=\textwidth}
\begin{tabular}{l r rrrrr}
\toprule
Model & Native & 48px & 96px & 112px & 160px & 224px \\
\midrule
R3D-18      & 112px & 56.4 & 72.6 & \textbf{71.5} & 74.9 & 76.1 \\
MC3-18      & 112px & 46.3 & 62.3 & \textbf{64.3} & 73.8 & 65.7 \\
R2+1D       & 112px & 61.6 & 73.1 & \textbf{76.8} & 79.4 & 76.2 \\
SlowFast    & 224px & 36.3 & 67.2 & 69.6 & 76.3 & \textbf{78.2} \\
\midrule
TimeSformer & 224px & 24.3 & 48.9 & 50.8 & 62.2 & \textbf{66.4} \\
ViViT       & 224px & 21.3 & 39.1 & 40.8 & 49.6 & \textbf{69.9} \\
VideoMAE    & 224px & 25.4 & 51.1 & 57.5 & 64.8 & \textbf{78.0} \\
VideoMamba  & 224px & 29.1 & 51.6 & 58.2 & 66.0 & \textbf{73.2} \\
\bottomrule
\end{tabular}
\end{adjustbox}
\end{table}

% ── S10: Dashboard and code availability ──────────────────────────────────────
\section{Interactive Dashboard and Code Availability}
\label{sec:sup_dashboard}

\paragraph{Interactive visualization dashboard.}
An anonymized interactive dashboard is included with the supplementary material for review. The public dashboard URL will be released after the review period. The dashboard provides:
\begin{enumerate}[leftmargin=*,topsep=2pt,itemsep=1pt]
\item \textbf{Temporal explorer}: coverage$\times$stride heatmaps per model and
  dataset, with configurable scaling and stride annotations.
\item \textbf{Spatial browser}: accuracy curves across resolutions for all
  8 models and datasets; target resolution tuning by family.
\item \textbf{Recommender}: \revdel{the system combines structured analysis with LLM synthesis through an external provider.} \rev{Given a target dataset, latency budget, frame rate, and deployment context, the review package applies deterministic rules from the paper tables to recommend a model--resolution--stride operating point. Any optional natural-language summarizer is disabled by default and is not used for any scientific result.}
\item \textbf{TDS comparison}: rank any subset of the 8 datasets by temporal demand,
  with model-specific breakdowns and Spearman correlation display.
\item \textbf{Routing simulator}: \rev{maximum-probability threshold sweep}\revdel{entropy threshold sweep} for any model on any dataset,
  with visual Pareto-frontier display.
\end{enumerate}

\begin{figure}[!htbp]
\centering
\begin{minipage}[t]{0.48\textwidth}
  \centering
  \includegraphics[width=\textwidth]{images/app1.png}
  \small{(a) Accuracy Explorer: native resolution (stride=16, cov=100\%, 224px). Summary table shows top-1 across all 8 datasets.}
\end{minipage}\hfill
\begin{minipage}[t]{0.48\textwidth}
  \centering
  \includegraphics[width=\textwidth]{images/app2.png}
  \small{(b) Accuracy Explorer: low resolution (stride=1, cov=75\%, 48px). Retrained CNN checkpoints are used at 48px.}
\end{minipage}\\[8pt]
\begin{minipage}[t]{0.48\textwidth}
  \centering
  \includegraphics[width=\textwidth]{images/app3.png}
  \small{\rev{(c) Architecture Recommender: deterministic FineGym example grounded only in the packaged result tables.}}
\end{minipage}\hfill
\begin{minipage}[t]{0.48\textwidth}
  \centering
  \includegraphics[width=\textwidth]{images/app4.png}
  \small{(d) Latency \& Efficiency: measured forward-pass latency (ms/clip, bfloat16, RTX PRO 6000) vs.\ spatial resolution for all 8 architectures. VideoMamba (green) is 4--8$\times$ slower than VideoMAE across all resolutions.}
\end{minipage}
\caption{Interactive dashboard screenshots from the anonymized supplementary package. The Accuracy Explorer (a,b) lets users sweep coverage, stride, and resolution interactively and view per-dataset top-1 accuracy tables. \rev{The Architecture Recommender (c) uses deterministic rules grounded in packaged measurements; no external API is required.} The Latency \& Efficiency tab (d) displays GPU latency curves across all model--resolution pairs.}
\label{fig:dashboard_screenshots}
\end{figure}

\paragraph{Code availability.}
Evaluation scripts, model fine-tuning code, result tables, and the Streamlit dashboard source are provided as an anonymized supplementary package for review. This includes the target-resolution pipeline with bicubic positional embedding interpolation. The public repository will be released after review.

% ── S11: Leave-one-architecture-out TDS stability ─────────────────────────────
\section{Leave-One-Architecture-Out \texorpdfstring{\TDS{}}{TDS} Stability}
\label{sec:sup_loo}

To assess whether the \TDS{} ranking depends on the specific 8-architecture pool evaluated, we recompute \TDS{} for all 8 datasets after removing each architecture in turn (7-architecture pools) and compare the resulting ranking to the full 8-architecture ranking via Spearman $\rho$.

\begin{table}[!htbp]
\centering
\caption{Leave-one-architecture-out \TDS{} ranking stability. $\rho$ is the Spearman correlation between the 8-dataset ranking computed with the named architecture excluded (7-architecture pool) and the full 8-architecture ranking.}
\label{tab:loo_tds}
\scriptsize\tabtight
\begin{tabular}{l r r}
\toprule
Excluded architecture & Spearman $\rho$ & $p$ \\
\midrule
R3D-18       & 1.000 & $<10^{-4}$ \\
MC3-18       & 1.000 & $<10^{-4}$ \\
R2+1D        & 1.000 & $<10^{-4}$ \\
SlowFast     & 1.000 & $<10^{-4}$ \\
TimeSformer  & 1.000 & $<10^{-4}$ \\
ViViT        & 1.000 & $<10^{-4}$ \\
VideoMAE     & 1.000 & $<10^{-4}$ \\
VideoMamba   & 0.976 & $3.3\times10^{-5}$ \\
\bottomrule
\end{tabular}
\end{table}

For 7 of 8 leave-one-out pools, the ranking is \emph{exactly} identical to the full-pool ranking ($\rho=1.00$). One leave-one-out pool swaps only the \#1/\#2 positions ($\rho=0.976$); ranks 3--8 are unaffected by every exclusion. \rev{This is sensitivity of the architecture-averaged estimator, not a protocol difference: every model--dataset pair is processed identically.} We therefore report FineGym and AUTSL as co-equal highest-\TDS{} domains rather than asserting a strict order.

\subsection{Architecture-Family Ablation}
\label{sec:sup_family_ablation}

Leave-one-out removes a single architecture at a time; a stronger test removes an entire family. We recompute \TDS{} under three restricted pools: \textbf{CNN-only} ($n{=}4$: R3D-18, MC3-18, R2+1D, SlowFast), \textbf{Transformer-only} ($n{=}3$: TimeSformer, ViViT, VideoMAE), and \textbf{Transformer+SSM} ($n{=}4$: adding VideoMamba). Table~\ref{tab:family_ablation} reports \TDS{} under each pool and Spearman $\rho$ against the full ranking.

\begin{table}[!htbp]
\centering
\caption{\TDS{} per dataset under architecture-family ablation pools, compared to the full 8-architecture ranking (sorted by full-pool \TDS{}).}
\label{tab:family_ablation}
\scriptsize\tabtight
\begin{tabular}{l r r r r}
\toprule
Dataset & Full ($n{=}8$) & CNN-only ($n{=}4$) & Transf.-only ($n{=}3$) & Transf.+SSM ($n{=}4$) \\
\midrule
FineGym       & 55.92 & 62.05 & 53.06 & 49.80 \\
AUTSL         & 53.02 & 67.20 & 46.45 & 38.83 \\
SSv2          & 27.29 & 32.16 & 25.95 & 22.43 \\
DriveAct      & 21.48 & 24.05 & 23.21 & 18.92 \\
Diving-48     & 19.16 & 21.74 & 20.59 & 16.59 \\
HMDB-51       & 16.47 & 21.47 & 14.54 & 11.47 \\
EPIC-Kitchens &  9.77 & 12.89 &  8.27 &  6.64 \\
UCF-101       &  4.86 &  7.36 &  2.90 &  2.35 \\
\midrule
Spearman $\rho$ vs.\ full & --- & 0.976 & 1.000 & 1.000 \\
\bottomrule
\end{tabular}
\end{table}

The ranking is essentially invariant to architecture-family composition: removing all four CNNs, keeping only the three Transformers, or keeping Transformers+SSM all leave ranks 3--8 in identical order, and only the CNN-only pool reproduces the same FineGym/AUTSL swap already characterized in the leave-one-out test above ($\rho=0.976$ in both cases, for the same reason: CNNs are comparatively more fragile on AUTSL's fine handshape transitions than on FineGym's phase boundaries, which narrows the gap). This directly addresses whether \TDS{} reflects a property of the dataset pool versus an artifact of including any one architecture family: under every family-restricted pool we tested, the ranking is the same dataset ordering reported in the main paper, with the sole exception of the already-disclosed \#1/\#2 ambiguity.

\subsection{Bootstrap Confidence Intervals}
\label{sec:sup_bootstrap}

To quantify ranking uncertainty directly rather than only testing discrete leave-one-out/family pools, we bootstrap over architectures: we resample 8 architectures with replacement ($B{=}10{,}000$ resamples) from the original pool, recompute \TDS{} for each dataset on each resample, and report the resulting 95\% percentile confidence interval in Table~\ref{tab:bootstrap_ci}.

\begin{table}[!htbp]
\centering
\caption{Bootstrap 95\% CI for \TDS{} per dataset (resampling 8 architectures with replacement, $B{=}10{,}000$).}
\label{tab:bootstrap_ci}
\scriptsize\tabtight
\begin{tabular}{l r r}
\toprule
Dataset & \TDS{} (bootstrap mean) & 95\% CI \\
\midrule
FineGym       & 55.99 & [47.24, 63.86] \\
AUTSL         & 53.13 & [36.02, 67.03] \\
SSv2          & 27.35 & [20.40, 33.92] \\
DriveAct      & 21.57 & [13.92, 29.35] \\
Diving-48     & 19.23 & [10.67, 28.37] \\
HMDB-51       & 16.54 & [9.53, 24.21] \\
EPIC-Kitchens &  9.80 & [6.25, 13.22] \\
UCF-101       &  4.88 & [2.18, 8.36] \\
\bottomrule
\end{tabular}
\end{table}

All eight 95\% CIs are disjoint from zero, indicating non-zero temporal demand at the architecture-pool level. AUTSL's CI is the widest (width 31.0pp) because its per-architecture drops are the most heterogeneous (TimeSformer and VideoMamba near-flat, CNNs collapse sharply), while EPIC-Kitchens and UCF-101 have the narrowest CIs, consistent with their low and relatively architecture-independent \TDS{}.

We additionally bootstrap the FineGym$-$AUTSL gap directly (same resampling procedure, computing $\TDS(\text{FineGym})-\TDS(\text{AUTSL})$ on each resample): mean $=+2.92$pp, 95\% CI $=[-4.58,\,+11.73]$, with FineGym exceeding AUTSL in 73.9\% of resamples. The confidence interval includes zero, which is the appropriate statistical test of the co-equal-ranking claim: the data are consistent with FineGym narrowly leading AUTSL in expectation, but do not support rejecting the null hypothesis of equal \TDS{} at the 95\% level. This is the quantitative basis, rather than only the leave-one-out swap, for reporting FineGym and AUTSL as co-equal highest-\TDS{} domains in the main paper.

% ── S12: Spectral and evidence-localization checks ───────────────────────────
\section{\rev{Spectral and Evidence-Localization Checks}}
\label{sec:sup_nyquist}

\rev{The main paper no longer interprets \TDS{} as literal signal aliasing: it measures accuracy degradation under the implemented stride and coverage protocol.} Here we report two diagnostics. The whole-frame optical-flow spectral proxy fails in the opposite direction from a naive signal-level Nyquist prediction; the second asks whether classifier evidence is temporally localized. Both are exploratory label-space diagnostics, not classical spectral validation.

\paragraph{Methodology.} For each of 8 datasets and 5 resolutions (48, 96, 112, 160, 224\,px), we sample up to 12 classes $\times$ 2 videos. We decode up to 64 \emph{consecutive} frames in native order, resize to the target resolution, and compute inter-frame optical-flow magnitude (Farneback) as a temporal signal sampled once per original frame. We estimate each signal's spectral cutoff frequency $f_c$ (90\% cumulative periodogram power, cycles/frame) and average over sampled videos to obtain one $f_c$ per (dataset, resolution) pair. We pair each $f_c$ with empirical stride-sensitivity at the same resolution: mean accuracy drop from stride=1 to stride=16 at coverage=100\%, averaged over 8 architectures and read from the existing sweep. This yields $8\times5=40$ (dataset, resolution) points.

\begin{table}[!htbp]
\centering
\caption{Mean flow-derived spectral cutoff frequency $f_c$, videos per resolution, and mean empirical stride-sensitivity, averaged over the 5 tested resolutions per dataset and sorted by stride-sensitivity.}
\label{tab:nyquist_v2}
\scriptsize\tabtight
\begin{tabular}{l r r r}
\toprule
Dataset & Mean $f_c$ & $n$/res. & Mean stride-sens. \\
\midrule
FineGym       & 0.188 & 24 & 46.30 \\
AUTSL         & 0.121 & 24 & 45.61 \\
SSv2          & 0.339 & 24 & 25.70 \\
DriveAct      & 0.192 & 21 & 18.65 \\
HMDB-51       & 0.383 & 24 & 16.57 \\
Diving-48     & 0.182 & 24 & 15.21 \\
EPIC-Kitchens & 0.248 & 23 &  9.70 \\
UCF-101       & 0.387 & 24 &  5.45 \\
\bottomrule
\end{tabular}
\end{table}

\paragraph{Result.} Pooling across all 40 (dataset, resolution) points, $f_c$ is significantly \emph{negatively} correlated with stride-sensitivity (Spearman $\rho=-0.635$, permutation $p=1.0\times10^{-5}$, $n=40$)---the opposite sign from the naive prediction that higher whole-frame temporal-frequency content should require denser sampling. FineGym integrates cleanly into this pattern: it has the highest mean stride-sensitivity (46.30pp) but a moderate whole-frame cutoff frequency ($f_c=0.188$ cycles/frame), close to DriveAct and Diving-48 rather than UCF-101/HMDB-51.

\paragraph{Interpretation.} \revdel{This is the same qualitative inversion already reported for FineGym in the main paper ($r=-0.475$), now shown across all 8 datasets with real statistical power rather than being a single-dataset anomaly.}\rev{The FineGym inversion is part of a broader descriptive pattern across the sampled dataset--resolution points.} UCF-101 and HMDB-51---appearance- and sports-driven domains with large whole-body and camera motion---have the highest flow-derived cutoff frequencies but among the lowest stride-sensitivity, because the action remains identifiable from almost any sampled frame. AUTSL and FineGym show the opposite behavior: their discriminative evidence is localized in handshape changes or brief gymnastics phase events, which are diluted by a whole-frame flow aggregate. \revdel{We conclude that whole-frame optical-flow frequency content is a systematically \emph{anti-correlated}, not merely noisy, proxy for the temporal demand that matters to a classifier.}\rev{In this sample, whole-frame optical-flow frequency is therefore an inverted proxy for classifier evidence demand, not a validation of a spectral mechanism.}

\subsection*{Label-Space Evidence-Localization Checks}

The inversion above arises because whole-frame optical flow measures \emph{pixel-level motion magnitude}, not the temporal localisation of \emph{discriminative evidence}. We conduct three complementary tests on FineGym ($n=97$ classes) using the fine-tuned TimeSformer checkpoint, which probe whether the model's evidence is temporally concentrated.

\paragraph{Test A: Temporal burstiness of optical flow (CPU).} Per-class, we compute the temporal coefficient of variation of inter-frame flow magnitudes ($\mathrm{CV} = \sigma_\mathrm{time}/\mu_\mathrm{time}$). After controlling for mean flow, temporal standard deviation positively predicts \rev{evidence loss}\revdel{aliasing}: $\partial r = +0.235$, $p=0.021$. The strong correlation between standard deviation and mean flow ($r=+0.909$) limits mechanistic interpretation.

\paragraph{Test B: TimeSformer temporal attention entropy (GPU).} We register forward hooks on the temporal self-attention Dropout layer in all 12 TimeSformer blocks, capturing the $[196\times12\times8\times8]$ temporal attention weight tensor per clip. Mean attention entropy positively correlates with \rev{evidence loss}\revdel{aliasing}: $r=+0.317$, $p=0.0015$. This is an association, not a mechanism test.

\paragraph{Test C: Sliding-window confidence concentration (GPU, most direct).} For each clip we decode 48 consecutive frames, slide a window of width $T=8$ with stride 4, and record correct-class probability. The concentration $1-H/\log n_\mathrm{windows}$ measures how localized the model's evidence is. Concentration positively predicts \rev{evidence loss}\revdel{aliasing}: $r=+0.263$, $p=0.009$; controlling for mean optical-flow magnitude gives $\partial r=+0.385$, $p<0.001$ ($n=97$ classes).

\paragraph{Summary.} Table~\ref{tab:label_space_nyquist} collects the three results.

\begin{table}[!htbp]
\centering
\caption{Label-space evidence-localization checks on FineGym ($n=97$ classes). All three tests use a fine-tuned TimeSformer trained on FineGym. Partial $r$ controls for whole-frame optical-flow magnitude.}
\label{tab:label_space_nyquist}
\scriptsize\tabtight
\begin{tabular}{l l r r r r}
\toprule
Test & Feature & $r$ & $p$ & partial $r$ & $p$ \\
\midrule
A: Optical flow burstiness & $\sigma_t(\text{flow}) / \mu_t(\text{flow})$ & $-0.421$ & ${<}0.001$ & $+0.235$ & $0.021$ \\
B: Temporal attention entropy & Mean attention $H$ (nats) & $+0.317$ & $0.0015$ & --- & --- \\
C: Sliding-window confidence & Confidence concentration & $+0.263$ & $0.009$ & $\mathbf{+0.385}$ & $\mathbf{<0.001}$ \\
\bottomrule
\end{tabular}
\end{table}

Test C provides the clearest classifier-level diagnostic: temporal concentration of correct-class confidence positively predicts per-class \rev{evidence loss}\revdel{aliasing}. Classes whose evidence is confined to a narrow window are more likely to lose that window under the protocol. This does not establish literal spectral aliasing.

\subsection*{Generalisation of Test C to All 8 Datasets}
\label{sec:sup_sw_generalisation}

We generalise Test C to all 8 datasets and an architecture ensemble. We correlate per-class confidence concentration with per-class \rev{evidence loss}\revdel{aliasing loss} (mean accuracy drop stride=1$\to$16) via Spearman $\rho$.

\paragraph{Window-geometry constraint.} The metric requires the decoded sequence to be substantially longer than the window $T$, so that windows are well separated and concentration reflects genuine temporal localisation. With 48 decoded frames, models with $T\le16$ satisfy this (neighbouring-window overlap $\le81\%$), but the two models with $T=32$ (ViViT, SlowFast) collapse to stride 1 and $97\%$ overlap: their 17 windows differ by a single frame, the confidence profile is near-flat, and the concentration measure floors out (mean concentration $0.10$ vs.\ $0.22$--$0.36$ for $T\le16$ models). On short-clip datasets such as FineGym (median ${\approx}37$ source frames) a 32-frame window cannot be slid at all. We therefore restrict the ensemble to the 6 architectures with $T\le16$ (TimeSformer, VideoMAE, R3D-18, MC3-18, R2+1D, VideoMamba); the $T=32$ models are a measurement artifact of the protocol rather than a property of the data.

\paragraph{Result.} Table~\ref{tab:sw_ensemble} reports positive correlations in all 8 datasets (pooled Spearman $\rho=+0.257$), significant in 4. \rev{We treat this as converging evidence for an association between temporal concentration and evidence loss, not as a causal or sampling-theoretic proof.}

\begin{table}[!htbp]
\centering
\caption{Sliding-window confidence-concentration test across all 8 datasets. Spearman $\rho$ relates per-class confidence concentration to per-class \rev{evidence loss}\revdel{aliasing loss}. The ensemble uses models with window $T\le16$; \textbf{bold} denotes $p<0.05$.}
\label{tab:sw_ensemble}
\scriptsize\tabtight
\begin{tabular}{l r r r r}
\toprule
Dataset & Classes & Models & Spearman $\rho$ & $p$ \\
\midrule
HMDB-51       &  51 & 6 & \textbf{+0.606} & ${<}0.001$ \\
UCF-101       & 101 & 6 & \textbf{+0.442} & ${<}0.001$ \\
DriveAct      &  33 & 6 & \textbf{+0.430} & 0.013 \\
SSv2          & 174 & 4 & \textbf{+0.269} & ${<}0.001$ \\
FineGym       &  97 & 6 & +0.138 & 0.177 \\
EPIC-Kitchens &  89 & 6 & +0.093 & 0.386 \\
AUTSL         & 226 & 6 & +0.074 & 0.270 \\
Diving-48     &  47 & 6 & +0.001 & 0.994 \\
\midrule
Pooled mean   &     &   & \textbf{+0.257} & --- \\
\bottomrule
\end{tabular}
\end{table}

Generalisation of Test B (TimeSformer temporal attention entropy) to the remaining datasets requires per-dataset attention extraction and is left to future work; the sliding-window concentration test (Test C), being architecture-agnostic, is the one we generalise here.

% ── S13: Additional robustness checks ───────────────────────────────────────
\section{Additional Robustness Checks}
\label{sec:sup_extra_checks}

\paragraph{Observational architecture interpretation.}
The architecture comparison in the main paper is intentionally phrased as an empirical association rather than a controlled causal ablation. Attention layout, state-space recurrence, tokenization, pretraining objective, parameter count, input length, native resolution, and optimization all covary across the evaluated architectures. A definitive mechanism test would require holding the backbone fixed while swapping only temporal aggregation modules. \revdel{Our contribution is the large controlled sampling sweep showing that, within this architecture pool, models with global divided attention or recurrent state integration are consistently less sensitive to temporal subsampling than CNNs and factorized-attention Transformers.}\rev{Moreover, the original sweep is confounded by native frame budget. Our revised architecture result therefore comes only from the matched-evidence comparison in Sec.~\ref{sec:sup_matched}: an accuracy-per-distinct-frame ordering remains, but neither causality nor the original 3--5$\times$ magnitude is supported.}

\paragraph{Dense-accuracy headroom does not explain \TDS{}.}
Because several datasets use balanced subsets or fixed 20-clips-per-class protocols, external leaderboard gaps are not directly comparable. We therefore compute internal dense-accuracy headroom. If under-optimization drove \rev{evidence loss}\revdel{aliasing loss}, headroom should correlate positively with stride=1$\to$16 drop. It does not: across 64 model--dataset pairs, Spearman $\rho=-0.22$ ($p=0.079$); at dataset level, mean headroom vs.\ \TDS{} gives $\rho=0.36$ ($p=0.385$).

\begin{table}[!htbp]
\centering
\caption{Dense-accuracy headroom check. Headroom is the mean gap from each model's stride=1 accuracy to the best stride=1 accuracy achieved by any model on the same dataset. It is not significantly correlated with \TDS{}, suggesting that temporal sensitivity is not an artifact of under-optimized dense accuracy.}
\label{tab:headroom_check}
\scriptsize\tabtight
\begin{tabular}{l r r r r}
\toprule
Dataset & Mean stride=1 & Best stride=1 & Mean headroom & \TDS{} \\
\midrule
FineGym       & 73.8 & 79.9 &  6.1 & 55.9 \\
AUTSL         & 72.9 & 82.3 &  9.3 & 53.0 \\
SSv2          & 42.5 & 52.4 &  9.9 & 27.6 \\
DriveAct      & 67.4 & 73.9 &  6.5 & 21.9 \\
Diving-48     & 40.3 & 53.0 & 12.7 & 19.2 \\
HMDB-51       & 78.2 & 84.0 &  5.8 & 16.6 \\
EPIC-Kitchens & 34.9 & 39.4 &  4.5 &  9.7 \\
UCF-101       & 89.0 & 95.5 &  6.5 &  4.9 \\
\bottomrule
\end{tabular}
\end{table}

\paragraph{Coverage--stride interaction.}
Figures in S1 visualize the complete coverage$\times$stride grids. To test additivity, we average accuracy over architectures for one high-\TDS{} and one low-\TDS{} dataset, then decompose each 5$\times$5 grid into row/column effects plus an interaction residual. FineGym has a meaningful interaction component: 15.0\% of grid variance, with RMSE 8.3pp. Low coverage is poor regardless of stride, while high coverage exposes the stride cliff. UCF-101 is closer to additive (4.0\%; RMSE 1.7pp).

\begin{table}[!htbp]
\centering
\caption{Mean accuracy (\%) over architectures for a high-\TDS{} dataset (FineGym) and a low-\TDS{} dataset (UCF-101). Rows are coverage; columns are stride.}
\label{tab:coverage_stride_examples}
\scriptsize\tabtight
\begin{adjustbox}{max width=\textwidth}
\begin{tabular}{l rrrrr @{\hskip18pt} l rrrrr}
\toprule
\multicolumn{6}{c}{FineGym} & \multicolumn{6}{c}{UCF-101} \\
\cmidrule(lr){1-6}\cmidrule(lr){7-12}
Coverage & s1 & s2 & s4 & s8 & s16 & Coverage & s1 & s2 & s4 & s8 & s16 \\
\midrule
10\%  &  8.0 &  7.2 &  6.4 &  5.8 &  5.8 & 10\%  & 74.5 & 73.2 & 69.5 & 64.2 & 57.7 \\
25\%  & 16.9 & 15.6 & 12.4 &  8.3 &  6.4 & 25\%  & 81.7 & 81.6 & 80.1 & 75.0 & 67.8 \\
50\%  & 42.9 & 40.7 & 28.7 & 17.1 &  9.3 & 50\%  & 86.6 & 86.6 & 85.7 & 82.9 & 76.6 \\
75\%  & 63.4 & 59.6 & 42.3 & 25.1 & 13.1 & 75\%  & 88.4 & 88.7 & 88.5 & 86.8 & 81.4 \\
100\% & 73.8 & 68.0 & 49.8 & 29.8 & 15.7 & 100\% & 89.0 & 89.0 & 89.0 & 88.1 & 84.1 \\
\bottomrule
\end{tabular}
\end{adjustbox}
\end{table}

\paragraph{\rev{Confidence-cascade calibration.}\revdel{Entropy-threshold calibration.}}
\revdel{The entropy router is not trained; it reuses two fixed inference outputs and sweeps a scalar confidence threshold on the evaluation split. The threshold should be selected on a held-out calibration split.} \rev{The cascade thresholds maximum class probability, not predictive entropy. We select $\tau$ on a stratified calibration half and score it on the untouched half over 1000 resamples. On SSv2/TimeSformer the held-out change relative to dense inference is $-0.03$pp (95\% CI $[-1.07,+0.83]$) at 6.9 average frames. Across 56 model--dataset pairs, however, the mean change is $-3.5$pp. The result supports confidence as a useful signal for some pairs, not a universally accuracy-neutral routing policy.}

% ── S14: What the stride axis varies ─────────────────────────────────────────
\section{\rev{What the Stride Axis Varies}}
\revnote
\color{revblue}
\label{sec:sup_padding}

Reviewers asked whether our temporal sweep measures aliasing or evidence loss.
This section documents the sampler's behaviour and the measurements that settle
the question.

\paragraph{Sampler.} For a clip of $T$ frames at coverage $\mathcal{C}$ and
stride $s$, with model frame budget $B$, the candidate pool is
$\{0,s,2s,\dots\}$ within the first $\lfloor T\mathcal{C}/100\rfloor$ frames.
If the pool holds at least $B$ entries it is re-uniformised to exactly $B$ by
linear interpolation over its index range; otherwise the shortfall is filled by
repeating the last candidate. Table~\ref{tab:sup_indices} makes both regimes
concrete.

\begin{table}[!htbp]
\centering
\caption{Frame positions actually supplied to a $B{=}16$ model at
$\mathcal{C}{=}100\%$. For a long clip the selected positions are nearly
independent of stride; for a short clip the input degenerates into a few
distinct frames followed by a static tail.}
\label{tab:sup_indices}
\scriptsize\tabtight
\begin{tabular}{l l r l}
\toprule
Clip & $s$ & Distinct & Positions (last six) \\
\midrule
$T{=}400$ & 1  & 16 & 266, 293, 319, 346, 372, 399 \\
$T{=}400$ & 4  & 16 & 264, 292, 316, 344, 368, 396 \\
$T{=}400$ & 16 & 16 & 256, 288, 304, 336, 352, 384 \\
\midrule
$T{=}60$  & 1  & 16 & 39, 43, 47, 51, 55, 59 \\
$T{=}60$  & 8  &  8 & 56, 56, 56, 56, 56, 56 \\
$T{=}60$  & 16 &  4 & 48, 48, 48, 48, 48, 48 \\
\bottomrule
\end{tabular}
\end{table}

\paragraph{Sampling interval in seconds.} Expressing the axis physically
(Table~\ref{tab:sup_seconds}) shows the same thing without reference to the
implementation: over most of its range the stride axis does not change the rate
at which the clip is sampled. Frame rates here are measured from the files, not
assumed; three datasets differ from the values used in our earlier duration
analysis (HMDB-51 is 30 not 25\,fps, AUTSL 30 not 25, EPIC-Kitchens 50 not 60).

\begin{table}[!htbp]
\centering
\caption{Median seconds between sampled frames at $\mathcal{C}{=}100\%$. At a
fixed nominal stride the realised interval differs by up to $3.9\times$ across
datasets, so equal-stride comparisons are not equal-rate comparisons.}
\label{tab:sup_seconds}
\scriptsize\tabtight
\begin{tabular}{l r rrrrr}
\toprule
Dataset & fps & $s{=}1$ & $s{=}2$ & $s{=}4$ & $s{=}8$ & $s{=}16$ \\
\midrule
UCF-101       & 25 & 0.395 & 0.395 & 0.395 & 0.395 & 0.617 \\
Diving-48     & 25 & 0.265 & 0.265 & 0.265 & 0.314 & 0.600 \\
SSv2          & 12 & 0.255 & 0.255 & 0.325 & 0.620 & 1.167 \\
DriveAct      & 15 & 0.183 & 0.183 & 0.256 & 0.500 & 1.000 \\
HMDB-51       & 30 & 0.167 & 0.167 & 0.167 & 0.260 & 0.504 \\
AUTSL         & 30 & 0.125 & 0.125 & 0.131 & 0.252 & 0.475 \\
EPIC-Kitchens & 50 & 0.120 & 0.120 & 0.120 & 0.158 & 0.301 \\
\bottomrule
\end{tabular}
\end{table}

\paragraph{The effect is confined to the shortfall regime.} Restricting to clips
whose candidate pool never falls below $B$ at $s{=}16$, the stride $1\to16$
accuracy difference is $-0.51$pp on average, against $+10.74$pp over all clips.
This holds in \textbf{21 of 21} model--dataset pairs with at least 20 such clips.
On AUTSL, SSv2 and DriveAct no clip escapes the shortfall regime at $s{=}16$, so
the entire reported effect on those datasets is measured inside it. The unpadded
subset is necessarily the longer clips, so this comparison is length-confounded
on its own; it should be read together with the index and interval tables above,
which do not depend on any subsetting.

\paragraph{Relation to \TDS{}.} Across datasets, the fraction of clips in the
shortfall regime correlates with \TDS{} at Spearman $0.964$. The relationship is
gating rather than explanatory: among AUTSL, SSv2 and DriveAct, all of which are
fully saturated, \TDS{} still spans $21.9$--$53.0$pp. Clip length relative to
frame budget determines whether temporal demand is expressed at all; dataset
content determines how large it is.

% ── S15: Matched-evidence protocol ───────────────────────────────────────────
\section{\rev{Matched-Evidence Architecture Comparison}}
\revnote
\label{sec:sup_matched}

Because frame budget $B$ governs when each architecture enters the shortfall
regime, and $B$ is rank-correlated with the reported drop at Spearman $0.849$
($p{=}0.008$), the uncontrolled comparison cannot isolate temporal aggregation
design. We therefore supply every architecture the same $k$ distinct frames at
the same positions, uniformly spanning the clip, resampled to each model's own
input length by uniform interpolation rather than end-repetition. Only
$k\leq8$ is strictly matched, since the smallest budget in the pool is 8.

Table~\ref{tab:sup_matched_models} reports the complete matched results over the
seven datasets with outputs for all eight architectures. Rankings are moderately
stable between adjacent budgets (Spearman $0.69$--$0.83$). At $k{=}2$, accuracy
agrees with the ordering under the original protocol (Spearman $0.857$,
$p{=}0.007$); the correlations at $k{=}1,4,8$ are $0.643$, $0.619$, and $0.476$
and are not significant. Ranked instead by the absolute accuracy drop between
$k{=}8$ and $k{=}2$, the ordering does not reproduce the original one (Spearman
$0.14$, $p{=}0.76$).

\begin{table}[!htbp]
\centering
\caption{Matched-evidence Top-1 accuracy (\%), averaged over the seven datasets
with complete outputs for all eight architectures. Every model receives the
same $k$ distinct frame positions. $\Delta_{8\to2}$ is the absolute accuracy
drop from $k=8$ to $k=2$.}
\label{tab:sup_matched_models}
\scriptsize\tabtight
\begin{tabular}{l rrrrr}
\toprule
Model & $k=1$ & $k=2$ & $k=4$ & $k=8$ & $\Delta_{8\to2}$ \\
\midrule
TimeSformer & 21.1 & 33.8 & 52.1 & 58.6 & 24.8 \\
VideoMamba  & 21.9 & 28.4 & 41.0 & 44.5 & 16.1 \\
MC3-18      & 24.7 & 29.5 & 39.8 & 51.0 & 21.6 \\
R3D-18      & 18.2 & 24.8 & 36.6 & 50.1 & 25.4 \\
R2+1D       & 18.0 & 25.3 & 37.5 & 50.5 & 25.3 \\
ViViT       & 20.4 & 27.6 & 41.3 & 53.6 & 25.9 \\
VideoMAE    & 17.5 & 23.2 & 38.2 & 48.7 & 25.5 \\
SlowFast    & 18.4 & 24.1 & 32.6 & 38.5 & 14.4 \\
\bottomrule
\end{tabular}
\end{table}

The budget confound is removed by construction. Consistent with frame budget
being a major confound,
the rank correlation between input length and measured degradation falls from
$0.849$ ($p{=}0.008$) under the published protocol to $0.231$ ($p{=}0.58$) under
matched evidence.
The spread in absolute drop narrows from $4.1\times$ under the original protocol
to $1.8\times$ under matched evidence. We therefore withdraw the original
magnitude and mechanism claims; the matched experiment supports only a
model-level evidence-efficiency comparison.

\rev{All eight architectures are included under the same protocol.}

% ── S16: Repeated-measures ANOVA ───────────────────────────
\section{\rev{Repeated-Measures Analysis of the Coverage\texorpdfstring{$\times$}{x}Stride Grid}}
\revnote
\label{sec:sup_rmanova}

All 25 cells of the grid score the same clips, so treating them as independent
observations, as our original two-way ANOVA did, understates the error term. We
re-estimate with clip as subject and coverage and stride as fully crossed
within-subject factors, testing each effect against its own subject$\times$factor
interaction and reporting partial $\eta^2$. Only clips present in all 25 cells
are used, so the design stays balanced.

Over 56 model--dataset pairs, partial $\eta^2$ is $0.292\pm0.152$ for coverage,
$0.232\pm0.145$ for stride and $0.070\pm0.081$ for their interaction. Coverage
therefore remains dominant but by $1.26\times$ rather than the $2.2\times$
implied by the between-cell analysis. The interaction, which the original model
omitted entirely, is largest on the high-demand datasets, reaching $0.271$ on
SlowFast/AUTSL, $0.262$ on R3D-18/AUTSL and $0.232$ on SlowFast/FineGym. The
per-model ordering of the stride effect is unchanged under the corrected model
(TimeSformer $0.124$, MC3-18 $0.197$, R3D-18 $0.226$, R(2+1)D $0.238$, ViViT
$0.247$, VideoMAE $0.259$, SlowFast $0.331$), so the qualitative grouping in the
main paper does not depend on the error term.

% ── S17: Held-out routing ────────────────────────────────────────────────────
\section{\rev{Held-Out Evaluation of the Confidence Cascade}}
\revnote
\label{sec:sup_routing_heldout}

The cascade thresholds the maximum class probability, not the predictive
entropy; our earlier name for it was inaccurate. More substantively, the
threshold was previously selected and scored on the same split. We repeat the
evaluation with $\tau$ chosen on a stratified calibration half and scored on the
untouched half, over $1000$ resamples.

On SSv2 with TimeSformer, the advantage over fixed dense inference falls from
$+0.21$pp in-sample to $-0.03$pp held-out, with a 95\% interval of
$[-1.07,+0.83]$: the cascade matches dense accuracy at $6.9$ average frames
rather than exceeding it. That cell is also the most favourable one. Averaged
over all 56 model--dataset pairs the held-out change relative to fixed dense
inference is $-3.5$pp, and 25 of 56 pairs have intervals excluding zero, with
losses reaching $-14.7$pp (ViViT/FineGym) and $-10.5$pp (VideoMAE/FineGym). The
supportable claim is that confidence carries usable routing signal on some
model--dataset pairs, not that temporal budgets can be halved without cost.

% ── S18: Convergence at low resolution ───────────────────────────────────────
\section{\rev{Convergence of Low-Resolution Fine-Tuning}}
\revnote
\label{sec:sup_convergence}

To check whether the 10-epoch budget limits the reported off-native accuracies,
we recovered the validation curves available at the five reported target resolutions
($n{=}503$ logs) and asked whether each run was still improving when the budget
ended.

\begin{table}[!htbp]
\centering
\caption{Convergence by target resolution. ``Still improving'' counts runs whose
best validation accuracy occurred at the final or penultimate epoch.}
\label{tab:sup_convergence}
\scriptsize\tabtight
\begin{tabular}{r r r r r}
\toprule
Resolution & Runs & Mean best epoch & Still improving & Gain in last 3 epochs \\
\midrule
48px  &  36 & 6.97 & 52.8\% & 1.38pp \\
96px  &  96 & 5.08 & 24.0\% & 4.21pp \\
112px &  93 & 4.94 & 16.1\% & 3.89pp \\
160px &  95 & 5.79 & 27.4\% & 3.27pp \\
224px & 183 & 5.66 & 23.0\% & 2.89pp \\
\bottomrule
\end{tabular}
\end{table}

The concern is partly borne out: 48px is the resolution where the budget binds
most, with over half the runs peaking at the final epoch against 16--27\%
elsewhere. The accuracy still being won in the closing epochs there is however
the smallest of any resolution, $1.38$pp, so the runs were converging slowly
rather than being cut off mid-ascent. We report the 48px numbers with this
limitation stated rather than claiming convergence.

% ── S19: TDS metric variants ─────────────────────────────────────────────────
\section{\rev{Sensitivity of the \texorpdfstring{\TDS{}}{TDS} Ranking to Its Definition}}
\revnote
\label{sec:sup_tds_variants}

\TDS{} was criticised for using only the endpoints of the stride curve, for
depending on baseline accuracy and floor effects, and for discarding negative
drops through its positive part. Each objection defines a variant, and none
changes the ranking (Table~\ref{tab:sup_tds_variants}).

\begin{table}[!htbp]
\centering
\caption{Dataset demand under five definitions, and the rank correlation of each
against the published \TDS{}. \emph{unclipped} removes $[\cdot]^+$;
\emph{AUC} integrates the whole stride curve trapezoidally in $\log_2 s$;
\emph{relative} divides by stride-1 accuracy; \emph{normalised} divides by the
headroom above chance.}
\label{tab:sup_tds_variants}
\scriptsize\tabtight
\begin{tabular}{l rrrrr}
\toprule
Dataset & Published & Unclipped & AUC & Relative & Normalised \\
\midrule
FineGym       & 55.92 & 55.92 & 24.33 & 76.92 & 78.03 \\
AUTSL         & 53.02 & 53.02 & 17.74 & 71.22 & 71.65 \\
SSv2          & 27.57 & 27.57 &  9.61 & 65.31 & 66.23 \\
DriveAct      & 21.88 & 21.88 &  8.84 & 31.87 & 33.35 \\
Diving-48     & 19.16 & 19.16 &  5.10 & 45.18 & 47.66 \\
HMDB-51       & 16.64 & 16.64 &  4.94 & 21.05 & 21.58 \\
EPIC-Kitchens &  9.74 &  9.74 &  3.02 & 26.86 & 27.73 \\
UCF-101       &  4.88 &  4.88 &  0.81 &  5.54 &  5.61 \\
\midrule
Spearman vs.\ published & --- & 1.000 & 1.000 & 0.952 & 0.952 \\
\bottomrule
\end{tabular}
\end{table}

The positive part never activates: no dataset has a negative mean drop, so
\emph{unclipped} is numerically identical to the published metric. Integrating
the full curve preserves the ordering exactly. Normalising by baseline accuracy
or by headroom above chance exchanges only adjacent pairs (DriveAct with
Diving-48, HMDB-51 with EPIC-Kitchens); the three highest-demand datasets are
identical under all five definitions.

% ── S20: Multi-stride augmentation ablation ──────────────────────────────────
\section{\rev{Multi-Stride Augmentation Ablation}}
\revnote
\label{sec:sup_tra}

A reviewer asked whether the measured cliffs reflect a mismatch between the
training and inference sampling distributions rather than anything intrinsic. We
tested this directly by fine-tuning with random $(\mathcal{C},s)$ exposure drawn
from the same grid used at evaluation, and comparing against the published
checkpoint. All three conditions are evaluated under the \emph{published}
sampler: evaluating the augmented arms without its padding would remove the
cliff by construction and answer nothing.

Two augmented arms are trained. \textbf{paper} draws its augmentation through the
published sampler, padding included; \textbf{fixed} draws from the same grid but
resamples uniformly, with no repeated final frame. Within each architecture the
two arms reach almost identical dense-sampling accuracy (TimeSformer 64.0\% and
63.6\%), so any difference in the stride curve is attributable to the form of the
augmentation rather than to one arm having trained better.

\begin{table}[!htbp]
\centering
\caption{Both architectures on AUTSL, evaluated under the published sampler at
coverage 100\%. ``Recovered'' is the share of the baseline stride-1 to stride-16
cliff regained at stride 16; ``dense cost'' is the accuracy given up at stride 1.}
\label{tab:sup_tra}
\scriptsize\tabtight
\begin{tabular}{l l rrrrr r r r}
\toprule
Model & Arm & $s{=}1$ & $s{=}2$ & $s{=}4$ & $s{=}8$ & $s{=}16$ & Cliff & Recov. & Dense cost \\
\midrule
\multirow{3}{*}{TimeSformer}
 & baseline & 66.8 & 66.9 & 66.3 & 65.7 & 50.8 & 16.0 & --- & --- \\
 & paper    & 63.5 & 63.6 & 63.1 & 62.6 & 55.3 &  8.2 & 28.4\% & $-3.3$ \\
 & fixed    & 63.1 & 63.4 & 63.4 & 62.3 & 52.6 & 10.4 & 11.7\% & $-3.7$ \\
\midrule
\multirow{3}{*}{R3D-18}
 & baseline & 75.0 & 75.3 & 68.0 & 27.2 &  6.5 & 68.5 & --- & --- \\
 & paper    & 65.1 & 65.3 & 62.1 & 44.9 & 31.9 & 33.2 & 37.0\% & $-9.9$ \\
 & fixed    & 63.8 & 63.2 & 59.9 & 28.0 & 10.4 & 53.3 &  5.7\% & $-11.2$ \\
\bottomrule
\end{tabular}
\end{table}

\rev{On AUTSL for these two architectures, multi-stride exposure materially reduces the measured cliff: the paper-sampler arm recovers 28.4\% for TimeSformer and 37.0\% for R3D-18. This establishes a train--inference distribution effect, but it is not a causal decomposition of the remaining loss. The augmented models retain 63--72\% of the baseline cliff and give up 3.3pp (TimeSformer) or 9.9pp (R3D-18) at stride 1, so the intervention trades dense and sparse operating points rather than producing universal robustness.}

\rev{The comparison between the two augmented arms is consistent with sensitivity to sampler details. Both see the same nominal $(\mathcal{C},s)$ grid and achieve similar dense accuracy, but the arm trained through the published repeat-last sampler recovers more at stride 16 than the fixed-length arm. This is consistent with adaptation to the train--inference input distribution. It does not, by itself, identify a spectral mechanism or prove that the residual loss is caused by missing temporal evidence. Broader model--dataset ablations would be needed for either conclusion.}

We therefore retain the conclusion that the dominant term is evidence loss, with
the correction that a recoverable mismatch term of 28--37\% sits alongside it.

% ── S21: The sampling-rate axis ──────────────────────────────────────────────
\section{\rev{Varying the Sampling Rate at Fixed Frame Count}}
\revnote
\label{sec:sup_span}

A reviewer observed that our temporal protocol investigates ``how early the
model has enough temporal context, rather than the aliasing due to frame
sampling stride''. This sweep tests that reading directly. The frame count $k$ is
pinned to each model's native budget and only the width $W$ of a centred window
varies, so the realised sampling rate $k/W$ changes by an order of magnitude
while the number of frames delivered never changes.

At fixed $k$, widening $W$ lowers the sampling rate but also admits more of the
action. These push in opposite directions, and no frame-subsampling design can
separate them --- that is the identifiability limitation stated in
Sec.~\ref{sec:intro}. What the sweep does decide is which one binds empirically.

\begin{table}[!htbp]
\centering
\caption{Top-1 (\%) against window width, averaged over architectures, with $k$
fixed at each model's frame budget. Widening the window reduces the sampling
rate tenfold, yet accuracy rises in seven of eight datasets.}
\label{tab:sup_span}
\scriptsize\tabtight
\begin{tabular}{l rrrrr r r}
\toprule
Dataset & $W{=}10\%$ & 25\% & 50\% & 75\% & 100\% & $\rho(W,\mathrm{acc})$ & $\Delta$ \\
\midrule
AUTSL         & 17.9 & 37.2 & 61.6 & 71.6 & 68.3 & $+1.000$ & $+50.4$ \\
SSv2          &  8.3 & 17.7 & 29.5 & 36.1 & 38.4 & $+1.000$ & $+30.1$ \\
Diving-48     & 12.8 & 22.3 & 38.3 & 43.2 & 41.5 & $+0.900$ & $+28.7$ \\
HMDB-51       & 56.5 & 68.0 & 76.5 & 80.4 & 80.2 & $+0.900$ & $+23.7$ \\
DriveAct      & 49.1 & 57.1 & 63.0 & 66.6 & 68.7 & $+1.000$ & $+19.6$ \\
EPIC-Kitchens & 17.4 & 24.8 & 31.5 & 35.2 & 36.4 & $+1.000$ & $+19.0$ \\
UCF-101       & 77.9 & 85.0 & 88.9 & 90.0 & 89.9 & $+0.900$ & $+12.0$ \\
\midrule
Kinetics-400  & 61.5 & 64.5 & 64.6 & 64.1 & 63.1 & $+0.100$ & $+1.6$ \\
\bottomrule
\end{tabular}
\end{table}

Seven of eight datasets are context-bound, four of them at $\rho=+1.000$: giving
a model the same number of frames spread over the whole action beats giving it
the same number of frames densely covering a tenth of it, by as much as $56$pp.
If undersampling were the dominant cost, the sign would be the other way round.
Over the range this protocol can reach, the binding constraint is how much of
the action is observed, not how densely it is sampled --- which is the reviewer's
reading, and we adopt it.

Kinetics-400 is the single exception, and it is flat rather than reversed
($\rho=+0.100$, $+1.6$pp across the whole range). This is interpretable: its
clips are fixed 10-second segments with the labelled action near the centre, so
the narrowest window already contains the discriminative evidence and widening it
mostly admits unrelated context. That one dataset sits at zero while the other
seven rise steeply is evidence that the measure discriminates rather than merely
trending positive.

We stress what this does \emph{not} establish. Because $W$ moves context and
rate together, a density effect cannot be ruled out, only shown to be dominated.
We therefore use this experiment to delimit the claim rather than to assert a
literal spectral mechanism.

\color{black}

\bibliographystyle{splncs04}
\bibliography{main}

\end{document}
